% GENERATED FILE. Edit canonical lessons, then run npm run generate:books.
% canonical-source-hash: fnv1a64:113b53d9d93c7f89
% canonical-lessons: KA-C262-torana, KA-C262-kunkuma, KA-C262-holige, KA-C262-prasada, KA-C262-abharana

\chapter{Festoon of mango leaves, Kumkum, Holige, Food offered to a god and then shared, Jewellery}
\label{ch:ka-festoon-of-mango-leaves-kumkum-holige-food-offered-to-a-god-and-then-shared-jewellery}
\hlchaptermodality{262}

\begin{chapteropening}
\textbf{By the end of this chapter:} \emph{I can say a festoon of mango leaves, kumkum, holige, food offered to a god and then shared and jewellery.}

Complete the last lesson of chapter 262: I can say a festoon of mango leaves, kumkum, holige, food offered to a god and then shared and jewellery.
\end{chapteropening}

\section[\texorpdfstring{\kn{ತೋರಣ} (tōraṇa)}{ತೋರಣ (tōraṇa)}]{\kn{ತೋರಣ} (tōraṇa) — a festoon of mango leaves (over a doorway)}

\label{lesson:KA-C262-torana}



\begin{quote}
Before the new one: say the Kannada for a lizard, then the Kannada for a cockroach.
\end{quote}

\subsection*{You'll want to know: \kn{ತೋರಣ}}
\textbf{\kn{ತೋರಣ}} — \emph{tōraṇa} — \textquotedblleft{}a festoon of mango leaves (over a doorway)\textquotedblright{}.

\begin{grammarlens}[title={What you've built}]
One more word for this chapter.
\end{grammarlens}

\subsection*{Your turn}
\begin{itemize}
\raggedright
  \item \emph{Say it:} \emph{tōraṇa}
  \item \emph{Say it:} \emph{tōraṇa}, once more
  \item \emph{Say it:} say \emph{tōraṇa}
  \item \emph{Recall:} say the Kannada for a lizard, then the Kannada for a cockroach, then say \emph{tōraṇa} again
\end{itemize}

\begin{tcolorbox}[breakable,colback=teal!4,colframe=teal!35!black,title={Before you move on}]
What does \kn{ತೋರಣ} mean? (\textquotedblleft{}A festoon of mango leaves (over a doorway)\textquotedblright{}.) Say it once more.
\end{tcolorbox}
\section[\texorpdfstring{\kn{ಕುಂಕುಮ} (kuṅkuma)}{ಕುಂಕುಮ (kuṅkuma)}]{\kn{ಕುಂಕುಮ} (kuṅkuma) — kumkum, red vermilion powder}

\label{lesson:KA-C262-kunkuma}



\begin{quote}
Before the new one: say the Kannada for a cockroach, then the Kannada for a festoon of mango leaves.
\end{quote}

\subsection*{You'll want to know: \kn{ಕುಂಕುಮ}}
\textbf{\kn{ಕುಂಕುಮ}} — \emph{kuṅkuma} — \textquotedblleft{}kumkum, red vermilion powder\textquotedblright{}.

\begin{grammarlens}[title={What you've built}]
One more word for this chapter.
\end{grammarlens}

\subsection*{Your turn}
\begin{itemize}
\raggedright
  \item \emph{Say it:} \emph{kuṅkuma}
  \item \emph{Say it:} \emph{kuṅkuma}, once more
  \item \emph{Say it:} say \emph{kuṅkuma}
  \item \emph{Recall:} say the Kannada for a cockroach, then the Kannada for a festoon of mango leaves, then say \emph{kuṅkuma} again
\end{itemize}

\begin{tcolorbox}[breakable,colback=teal!4,colframe=teal!35!black,title={Before you move on}]
What does \kn{ಕುಂಕುಮ} mean? (\textquotedblleft{}Kumkum, red vermilion powder\textquotedblright{}.) Say it once more.
\end{tcolorbox}
\section[\texorpdfstring{\kn{ಹೋಳಿಗೆ} (hōḷige)}{ಹೋಳಿಗೆ (hōḷige)}]{\kn{ಹೋಳಿಗೆ} (hōḷige) — holige, a sweet stuffed flatbread}

\label{lesson:KA-C262-holige}



\begin{quote}
Before the new one: say the Kannada for a festoon of mango leaves, then the Kannada for kumkum.
\end{quote}

\subsection*{You'll want to know: \kn{ಹೋಳಿಗೆ}}
\textbf{\kn{ಹೋಳಿಗೆ}} — \emph{hōḷige} — \textquotedblleft{}holige, a sweet stuffed flatbread\textquotedblright{}.

\begin{grammarlens}[title={What you've built}]
One more word for this chapter.
\end{grammarlens}

\subsection*{Your turn}
\begin{itemize}
\raggedright
  \item \emph{Say it:} \emph{hōḷige}
  \item \emph{Say it:} \emph{hōḷige}, once more
  \item \emph{Say it:} say \emph{hōḷige}
  \item \emph{Recall:} say the Kannada for a festoon of mango leaves, then the Kannada for kumkum, then say \emph{hōḷige} again
\end{itemize}

\begin{tcolorbox}[breakable,colback=teal!4,colframe=teal!35!black,title={Before you move on}]
What does \kn{ಹೋಳಿಗೆ} mean? (\textquotedblleft{}Holige, a sweet stuffed flatbread\textquotedblright{}.) Say it once more.
\end{tcolorbox}
\section[\texorpdfstring{\kn{ಪ್ರಸಾದ} (prasāda)}{ಪ್ರಸಾದ (prasāda)}]{\kn{ಪ್ರಸಾದ} (prasāda) — food offered to a god and then shared}

\label{lesson:KA-C262-prasada}



\begin{quote}
Before the new one: say the Kannada for kumkum, then the Kannada for holige.
\end{quote}

\subsection*{You'll want to know: \kn{ಪ್ರಸಾದ}}
\textbf{\kn{ಪ್ರಸಾದ}} — \emph{prasāda} — \textquotedblleft{}food offered to a god and then shared\textquotedblright{}.

\begin{grammarlens}[title={What you've built}]
One more word for this chapter.
\end{grammarlens}

\subsection*{Your turn}
\begin{itemize}
\raggedright
  \item \emph{Say it:} \emph{prasāda}
  \item \emph{Say it:} \emph{prasāda}, once more
  \item \emph{Say it:} say \emph{prasāda}
  \item \emph{Recall:} say the Kannada for kumkum, then the Kannada for holige, then say \emph{prasāda} again
\end{itemize}

\begin{tcolorbox}[breakable,colback=teal!4,colframe=teal!35!black,title={Before you move on}]
What does \kn{ಪ್ರಸಾದ} mean? (\textquotedblleft{}Food offered to a god and then shared\textquotedblright{}.) Say it once more.
\end{tcolorbox}
\section[\texorpdfstring{\kn{ಆಭರಣ} (ābharaṇa)}{ಆಭರಣ (ābharaṇa)}]{\kn{ಆಭರಣ} (ābharaṇa) — jewellery}

\label{lesson:KA-C262-abharana}



\begin{quote}
Before the new one: say the Kannada for holige, then the Kannada for food offered to a god and then shared.
\end{quote}

\subsection*{You'll want to know: \kn{ಆಭರಣ}}
\textbf{\kn{ಆಭರಣ}} — \emph{ābharaṇa} — \textquotedblleft{}jewellery\textquotedblright{}.

\begin{grammarlens}[title={What you've built}]
One more word for this chapter.
\end{grammarlens}

\subsection*{Your turn}
\begin{itemize}
\raggedright
  \item \emph{Say it:} \emph{ābharaṇa}
  \item \emph{Say it:} \emph{ābharaṇa}, once more
  \item \emph{Say it:} say \emph{ābharaṇa}
  \item \emph{Recall:} say the Kannada for holige, then the Kannada for food offered to a god and then shared, then say \emph{ābharaṇa} again
\end{itemize}

\begin{tcolorbox}[breakable,colback=teal!4,colframe=teal!35!black,title={Before you move on}]
What does \kn{ಆಭರಣ} mean? (\textquotedblleft{}Jewellery\textquotedblright{}.) Say it once more.
\end{tcolorbox}
